\documentclass[11pt]{article}
\usepackage[margin=1in]{geometry}
\usepackage{amsmath,amssymb,amsthm}
\usepackage{booktabs}
\usepackage{hyperref}

\newtheorem{assumption}{Assumption}
\newtheorem{proposition}{Proposition}
\newtheorem{remark}{Remark}

\newcommand{\E}{\mathbb{E}}
\newcommand{\Var}{\mathrm{Var}}
\newcommand{\Prob}{\Pr}
\newcommand{\ind}{\mathbf{1}}

\title{A Task-Level Cost Model for LLM Endpoint Selection\\
\large with prompt caching, failures, time value, and uncertainty penalties}
\author{Matthew Pollard}
\date{\today}

\begin{document}
\maketitle

\begin{abstract}
The price of a language-model call is not its list price. Over a multi-turn task
the transcript is resubmitted at every turn, so the caller pays for a growing
prefix $N$ times. Whether that prefix is served from a prompt cache, whether the
request fails and must be retried cold, and how long the caller waits all change
the bill, and all three differ between providers serving the same model. This
document defines the cost of a \emph{task} as the sum of its turn costs, derives
its expectation and a variance bound in closed form, and treats OpenRouter's
published 24-hour rates as measured inputs rather than parameters to be estimated.
All quantities map to fields published by OpenRouter or to explicit caller
assumptions.
\end{abstract}

\section{Setup}

\subsection{Task profile}

A task is a sequence of $N$ turns indexed $k = 1, \dots, N$. At each turn the
caller appends $a$ new prompt tokens (user text or tool results) and the model
emits $o$ completion tokens. The transcript therefore grows by
\[
  d = a + o
\]
tokens per turn. At turn $k$ the prompt is the full transcript so far. It
consists of a \emph{reusable prefix} that was already sent on the previous turn,
\[
  S_k = (k-1)\,d, \qquad k = 1,\dots,N,
\]
followed by the $a$ new tokens, which no cache can have seen. The total prompt at
turn $k$ is $S_k + a$. Inter-turn gaps are assumed short relative to cache
retention; see Section~\ref{sec:beyond} for the relaxation.

\subsection{Endpoint prices}

An endpoint $e$ publishes per-token prices, in USD per token (the tools display
them per million):
\begin{center}
\begin{tabular}{@{}lll@{}}
\toprule
Symbol & Meaning & OpenRouter field \\
\midrule
$w_e$ & uncached prompt token & \texttt{input\_cache\_write} if published, else \texttt{prompt} \\
$r_e$ & cached prompt token & \texttt{input\_cache\_read} \\
$c_e$ & completion token & \texttt{completion} \\
$f_e$ & fixed fee per request (USD) & \texttt{request} \\
\bottomrule
\end{tabular}
\end{center}
If \texttt{input\_cache\_read} is absent the endpoint has no cache and we set
$r_e = w_e$ and $h_e = 0$ below. Prices from the endpoints API are already net of
any advertised discount; the \texttt{discount} field is informational.

\subsection{Endpoint telemetry}

For each endpoint OpenRouter publishes 24-hour aggregates:
\begin{center}
\begin{tabular}{@{}lll@{}}
\toprule
Symbol & Meaning & Source \\
\midrule
$h_e \in [0,1]$ & token-weighted prompt cache hit rate & \texttt{cacheHitRate} \\
$T_e$ & tokens served in the window & \texttt{totalTokens} \\
$u_e \in [0,1]$ & uptime (fraction of successful requests) & \texttt{uptime\_last\_1d} \\
$\ell_e^{50}, \ell_e^{90}$ & p50 / p90 time to first token (s) & \texttt{p50\_latency}, \texttt{p90\_latency} \\
$g_e^{50}, g_e^{90}$ & p50 / p90 output throughput (tokens/s) & \texttt{p50\_throughput}, \texttt{p90\_throughput} \\
\bottomrule
\end{tabular}
\end{center}

\subsection{Caller parameters}
\begin{center}
\begin{tabular}{@{}lll@{}}
\toprule
Symbol & Meaning & Default \\
\midrule
$N, a, o$ & task profile & profile-dependent \\
$v$ & value of the caller's time, USD per hour & $0$ \\
$\kappa_\ell$ & latency pessimism, blend toward p90 & $0$ \\
$\lambda$ & risk aversion, in task standard deviations & $0$ \\
$B$ & fallback endpoint on failure & the endpoint itself, cold \\
\bottomrule
\end{tabular}
\end{center}

\section{Assumptions}

\begin{assumption}[Independent cache outcomes]\label{as:indep}
Let $H_k \in [0,1]$ be the fraction of the reusable prefix $S_k$ served from cache
at turn $k$. The $H_k$ are independent across turns with common mean
$\E[H_k] = h$, the endpoint's published hit rate.
\end{assumption}

The Bernoulli case, $H_k \in \{0,1\}$ with $\Prob[H_k = 0] = p = 1 - h$, is the
special case in which every turn is a complete hit or a complete miss. The
published $h$ is token-weighted and identifies the mean of $H_k$ but not its
distribution; the Bernoulli model is the maximum-variance member of the family,
which is used in Section~\ref{sec:var}.

\begin{assumption}[Independent failures]\label{as:fail}
Each request fails with probability $q = 1 - u$, independently of everything
else. A failed request is not billed for tokens. It is retried once on the
fallback endpoint $B$, which succeeds and has no cached prefix for this task.
\end{assumption}

\begin{assumption}[Time is linear in tokens]\label{as:time}
Waiting for the first token costs $\ell$ seconds and generating $o$ tokens
costs $o / g$ seconds. Both are valued at $v$ USD per hour. Time to first token
is treated as independent of prompt length because OpenRouter publishes only a
single p50 per endpoint.
\end{assumption}

\begin{assumption}[Missing telemetry is worst-case]\label{as:missing}
If an endpoint lacks a latency, throughput, or uptime observation, the missing
value is imputed as the worst observed value among endpoints serving the same
model. An endpoint is never made cheaper by having less data.
\end{assumption}

\section{Inputs are measurements, not estimates}\label{sec:inputs}

The hit rate $h$ and uptime $u$ enter the model exactly as OpenRouter publishes
them. Each is a token- or request-weighted average over every call the endpoint
served in the last 24 hours. For GLM 5.3 Flash on 2026-09-04 the smallest primary
endpoint served $3 \times 10^8$ tokens and the largest $4 \times 10^{10}$; even
at $10^4$ tokens per request that is tens of thousands to millions of requests,
and the sampling standard deviation of a proportion at those counts is below
$0.3$ percentage points. No shrinkage toward a prior is applied: any prior heavy
enough to move the estimate would dwarf the sampling noise and merely penalise
smaller providers for being smaller.

Two sources of uncertainty remain and are outside the model. The 24-hour window
moves, so the published rate drifts from day to day; a penalty for that should be
derived from the observed temporal variance of the rate, which requires recording
history. And the aggregate describes other callers' traffic, not the caller's own
task; the right instrument for that is an explicit override of $h$ (for example
$h = 0$ for a cold-cache bound), not a prior.

\subsection{Latency}

OpenRouter publishes p50 and p90 time to first token. A pessimistic latency is a
blend toward the tail,
\begin{equation}\label{eq:pess-l}
  \ell = (1 - \kappa_\ell)\, \ell^{50} + \kappa_\ell\, \ell^{90},
  \qquad \kappa_\ell \in [0, 1],
\end{equation}
and throughput uses the median, $g = g^{50}$, since the published p90 is the
favourable tail. If \texttt{input\_cache\_read} is absent the endpoint has no
cache and $h = 0$; if no hit rate is published for a cacheable endpoint it is
scored as cold, $h = 0$, consistent with Assumption~\ref{as:missing}.

\section{Turn cost}

Fix an endpoint $A$ with inputs $h, u, \ell, g$ and prices $w, r, c, f$. Write $v' = v / 3600$ for the value of one
second.

\subsection{Successful turn}

The $a$ new tokens are always billed at $w$. The prefix $S_k$ is billed at $r$
for the cached fraction and $w$ for the rest:
\begin{equation}\label{eq:turn-ok}
  X_k^{\text{ok}}
  = \underbrace{w\,a + c\,o + f}_{\text{fixed}}
  + \underbrace{v'\big(\ell + o / g\big)}_{\text{time}}
  + \underbrace{S_k\,\big[\, w - H_k\,(w - r)\,\big]}_{\text{prefix}} .
\end{equation}
Taking expectations under Assumption~\ref{as:indep},
\begin{equation}\label{eq:turn-ok-mean}
  \E\big[X_k^{\text{ok}}\big] = \alpha + \pi\, S_k,
  \qquad
  \alpha = w\,a + c\,o + f + v'\big(\ell + o/g\big),
  \qquad
  \pi = h\, r + (1 - h)\, w .
\end{equation}
$\pi$ is the \emph{effective prefix price}. The new tokens never receive it.

\subsection{Failed turn}

With probability $q$ the request to $A$ fails. The caller has waited
$\ell_A$ seconds for nothing, then repeats the turn on the fallback
$B$, whose cache holds none of this task's prefix. The retried turn costs
\begin{equation}\label{eq:turn-fail}
  X_k^{\text{fail}}
  = v' \ell_A
  + \underbrace{w_B\,(S_k + a) + c_B\, o + f_B + v'\big(\ell_B + o / g_B\big)}_{\text{cold turn on } B}
  + \underbrace{h\,(w - r)\, d}_{\text{return penalty}} .
\end{equation}
The return penalty is the expected extra cost at turn $k+1$: $A$'s cache
never saw turn $k$'s completion, so of the prefix $S_{k+1}$ only $S_k$ can be
cached and the remaining $d$ tokens are billed at $w$ rather than at the
blended price. Writing $\alpha_B = w_B\, a + c_B\, o + f_B + v'(\ell_B
+ o / g_B)$,
\begin{equation}\label{eq:turn-fail-mean}
  \E\big[X_k^{\text{fail}}\big] = \beta + w_B\, S_k,
  \qquad
  \beta = v' \ell_A + \alpha_B + h\,(w - r)\, d .
\end{equation}
When $B = A$ (the default: the endpoint retried cold) this reduces to the
prefix loss $h\,(w - r)\,S_k$ plus one wasted wait, which is the
failure term in the current scorer, plus the retry's own time.

\subsection{Expected turn cost}

Combining \eqref{eq:turn-ok-mean} and \eqref{eq:turn-fail-mean},
\begin{equation}\label{eq:turn-mean}
  \E[X_k] = (1 - q)\,(\alpha + \pi\, S_k) + q\,(\beta + w_B\, S_k)
  = \underbrace{(1-q)\,\alpha + q\,\beta}_{A_0}
  + \underbrace{\big[(1-q)\,\pi + q\, w_B\big]}_{\Pi}\, S_k .
\end{equation}
The turn cost is affine in the prefix length with intercept $A_0$ and slope
$\Pi$, the failure-adjusted effective prefix price.

\section{Task cost}

\begin{proposition}[Task cost]\label{prop:task}
The task cost is the sum of the turn costs, $X_{\text{task}} = \sum_{k=1}^N X_k$.
Under Assumptions~\ref{as:indep}--\ref{as:time} and constant growth
$S_k = (k-1)d$,
\begin{equation}\label{eq:task-mean}
  \boxed{\;
  \E[X_{\text{task}}]
  = N\, A_0 \;+\; \Pi\, d\, \frac{N(N-1)}{2}
  \;}
\end{equation}
\end{proposition}
\begin{proof}
Sum \eqref{eq:turn-mean} over $k$ and use $\sum_{k=1}^{N} (k-1) = N(N-1)/2$.
\end{proof}

\subsection{Decomposition}

Expanding $\Pi$ separates what the caller pays under a perfect cache from what
misses and failures add:
\begin{equation}\label{eq:decomp}
  \Pi\, d\, \frac{N(N-1)}{2}
  = \underbrace{r\, d\, \tfrac{N(N-1)}{2}}_{\text{cached-read baseline}}
  + \underbrace{(1-q)\,(1-h)\,(w - r)\, d\, \tfrac{N(N-1)}{2}}_{\text{miss premium}}
  + \underbrace{q\,(w_B - r)\, d\, \tfrac{N(N-1)}{2}}_{\text{failure premium}} .
\end{equation}
All three terms are $\Theta(N^2)$. There are $N$ submissions and the $k$-th
exposes a prefix of length $(k-1)d$, so every per-token surcharge on the
prefix is paid on a quadratically growing volume. A hit rate of $95\%$ is
therefore not a $5\%$ surcharge on a long task: the miss premium is
$\Theta((1 - h)\, N^2)$ while the linear term $N A_0$ does not grow with
context at all.

\begin{remark}[Mean-prefix equivalence]
Because \eqref{eq:turn-mean} is affine in $S_k$,
\[
  \E[X_{\text{task}}] = N \cdot \E\big[X_k \,\big|\, S_k = \bar S\big],
  \qquad \bar S = d\, \frac{N-1}{2}.
\]
The task equals $N$ representative turns at the mean prefix, not one. The
factor $N$ does not affect the ranking of endpoints but does affect totals.
\end{remark}

\begin{remark}[Relation to the single-turn scorer]
The current scorer takes one prompt size $C$ and applies $\pi$ to all of it.
That treats the new tokens as partly cached and underprices each turn by
$h\,(w - r)\, a$. It also reports one turn, not $N$. Replacing $C$ by
the pair $(a, S_k)$ and summing over $k$ recovers \eqref{eq:task-mean}.
\end{remark}

\subsection{Comparing endpoints}

For endpoints $A$ and $A'$ serving the same model,
\begin{equation}\label{eq:compare}
  \E[X^{A}_{\text{task}}] - \E[X^{A'}_{\text{task}}]
  = N\,\big(A_0 - A_0'\big) + \big(\Pi - \Pi'\big)\, d\, \frac{N(N-1)}{2}.
\end{equation}
When list prices coincide, as they do for every primary endpoint of many
open-weight models, the first term vanishes and the gap is purely quadratic
and driven by $h$ and $u$ alone.

\section{Variance and risk}\label{sec:var}

Under Assumption~\ref{as:indep} and ignoring the randomness of failures, the
only random terms are $H_k$, and
\begin{equation}\label{eq:var}
  \Var[X_{\text{task}}]
  = (w - r)^2\, d^2 \sum_{k=1}^{N} (k-1)^2\, \Var[H_k]
  \;\le\; (w - r)^2\, h\,(1 - h)\, d^2\, \frac{(N-1)\,N\,(2N-1)}{6}.
\end{equation}
The bound is attained by the Bernoulli model and is the largest variance of any
$[0,1]$-valued $H_k$ with mean $h$. The standard deviation grows as
$N^{3/2}$; the coefficient of variation shrinks as $N^{-1/2}$. The published
$h$ alone gives an upper bound on tail risk, not an estimate of it. The single
worst event is a full miss on the last turn, costing $(w - r)\, d\, (N-1)$.

Including failures, $X_k$ is a mixture of two affine functions of $S_k$ and
its variance adds $q(1-q)\big[(\beta - \alpha) + (w_B - \pi) S_k\big]^2$ per
turn; the closed form is routine and omitted.

\subsection{Decision rule}

Given a quality constraint $Q_m \ge q_{\min}$ on the model $m$ served, choose
\begin{equation}\label{eq:decision}
  \min_{(m, e)} \;\; \E\big[X_{\text{task}}(m, e)\big] + \lambda\, \sqrt{\Var\big[X_{\text{task}}(m, e)\big]}
  \quad \text{s.t.} \quad Q_m \ge q_{\min},
\end{equation}
with $\lambda = 0$ for a risk-neutral caller. The Pareto tools report the
non-dominated set over the objectives
\[
  \big(\E[X_{\text{task}}] \downarrow,\ \ell \downarrow,\ g \uparrow,\ h \uparrow,\ u \uparrow\big)
\]
so that a caller who will not commit to $v$ and $\lambda$ can still see the
trade-offs.

\section{Worked example}

GLM 5.3 Flash, primary endpoints, 2026-09-04. Every endpoint lists
$w = 0.075$, $r = 0.015$, $c = 0.25$ USD per million tokens, so only telemetry
separates them. Task profile $a = 2000$, $o = 500$, $N = 200$, $v = 0$, $B = A$.

\begin{center}
\begin{tabular}{@{}lrrrr@{}}
\toprule
Endpoint & $h$ & $u$ & $\Pi$ (\$/M) & $\E[X_{\text{task}}]$ \\
\midrule
Z.ai      & 0.859 & 0.998 & 0.0236 & \$1.23 \\
DeepInfra & 0.497 & 0.993 & 0.0454 & \$2.31 \\
\bottomrule
\end{tabular}
\end{center}

With $\sum_k S_k = d\,N(N-1)/2 = 49.75$M prefix tokens, the quadratic term is
\$1.17 on Z.ai and \$2.26 on DeepInfra, against a linear term of \$0.055 on
both. At the default single turn of the current tools the same two endpoints
differ by $25\%$; over the task they differ by $88\%$, and all of it comes from
the hit rate. The Bernoulli variance bound \eqref{eq:var} gives a standard
deviation of \$0.085 on Z.ai and \$0.122 on DeepInfra, so at $N = 200$ the two
endpoints are separated by roughly nine standard deviations of the noisier one.

\section{Beyond independent misses}\label{sec:beyond}

Assumption~\ref{as:indep} is the weakest link. Three refinements are stated as
explicit knobs because OpenRouter's data cannot identify them.

\paragraph{Eviction.}
Caches expire after a time-to-live $T_{\text{ttl}}$. With inter-turn gap
$\tau_k$, replace the constant hit rate by
$h_k = h \cdot \ind[\tau_k < T_{\text{ttl}}]$ or by
$h_k = h\, e^{-\tau_k / T_{\text{ttl}}}$. Tool-call bursts then hit and
human think-time misses.

\paragraph{Cache cap.}
If a provider caches at most $L$ tokens, the cached part of the prefix is
$\min(S_k, L)$ and the remainder is billed at $w$. Whether exceeding the cap
instead invalidates the whole prefix is provider-specific and should be a
switch, not an assumption.

\paragraph{Provider state.}
Let $Z_k$ be the endpoint serving turn $k$ under the routing policy and
failures, and give each endpoint its own cache state (longest held prefix and
time last touched). Turn $k$ hits on $Z_k$ only if that endpoint holds the
current prefix within its TTL. A return to $A$ after a detour to $B$ hits if
$A$'s copy is still alive. This is a Markov chain over (endpoint, cache
state); the independent model is its memoryless approximation.

\paragraph{Compounding.}
If a disruption at turn $t_j$ leaves the task cold for $R_j$ turns, the
premium is $(w - r) \sum_j \sum_{k = t_j}^{t_j + R_j - 1} S_k$. With fixed
$R_j$ this is still $O(N^2)$; if recovery scales with task length it is
$O(N^3)$; if a switch permanently loses affinity it is a regime change and the
cold cost applies from then on.

\section{Reporting}

The unit of optimisation is the task. For each endpoint the tools should report
\begin{itemize}
  \item $\E[X_{\text{task}}]$ in dollars per completed task, with the
        decomposition \eqref{eq:decomp}: fixed, time, cached-read baseline,
        miss premium, failure premium;
  \item the perfect-cache cost ($h = 1$) and cold-cache cost
        ($h = 0$) as bounds;
  \item the variance bound \eqref{eq:var} and the risk-adjusted objective
        \eqref{eq:decision};
  \item total fresh, cached-read, and completion tokens;
  \item the telemetry inputs $h, u, \ell, g$ and the timestamp of the
        24-hour window they summarise.
\end{itemize}
Per-million normalisation is secondary: repeatedly submitted transcript tokens
make the denominator ambiguous, and the figure falls with context while the
bill rises.

\end{document}
